\documentclass[10pt,twocolumn]{article}

\usepackage[utf8]{inputenc}
\usepackage[T1]{fontenc}
\usepackage[brazil]{babel}
\usepackage[a4paper,margin=1.8cm]{geometry}
\usepackage{times}
\usepackage{url}
\usepackage{enumitem}
\usepackage{balance}
\usepackage[backend=biber,style=abnt]{biblatex}
\usepackage{hyperref}

\addbibresource{relatorio-final-abnt.bib}

\hypersetup{colorlinks=true,linkcolor=black,urlcolor=blue,citecolor=black}
\setlength{\columnsep}{0.7cm}
\setlength{\parindent}{1.1em}
\setlength{\parskip}{0.15em}
\setlist{nosep,leftmargin=1.2em}

\title{\vspace{-1.2em}\textbf{Aurora: Quem Decide?}\\[0.2em]
\large Relat\'orio Final de um \emph{Serious Game} 2D de Educa\c{c}\~ao Pol\'itica e Cidadania}
\author{Projeto Acad\^emico --- MAC5784: Intelig\^encia Artificial em Jogos de Computador\\
Universidade de S\~ao Paulo (USP) --- 2026}
\date{}

\begin{document}
\maketitle

\begin{abstract}
\noindent \emph{Aurora: Quem Decide?} \'e um \emph{serious game} 2D \emph{top-down},
desenvolvido em Python com a biblioteca Arcade, cujo prop\'osito \'e ensinar de
forma interativa as atribui\c{c}\~oes e os limites dos cargos eletivos brasileiros
e a divis\~ao de compet\^encias entre os n\'iveis de governo. Este relat\'orio
apresenta os objetivos, o processo de desenvolvimento e os resultados obtidos
com o prot\'otipo (\emph{vertical slice}) e com a aplica\c{c}\~ao de apresenta\c{c}\~ao
que documenta as decis\~oes de projeto. O conte\'udo \'e educativo, fict\'icio e
n\~ao partid\'ario.
\end{abstract}

\section{Objetivos}

O objetivo pedag\'ogico central \'e ensinar, por meio da jogabilidade, \emph{quem
decide o qu\^e} na estrutura pol\'itico-administrativa brasileira --- as atribui\c{c}\~oes
de vereador, prefeito, deputado estadual, governador, deputado federal, senador e
presidente ---, abordando ainda or\c{c}amento p\'ublico, fiscaliza\c{c}\~ao, participa\c{c}\~ao
social e coopera\c{c}\~ao entre os n\'iveis de governo. Ao final de cada cap\'itulo, o
jogador deve ser capaz de: (i) \textbf{relacionar} um problema ao n\'ivel de governo
respons\'avel; (ii) \textbf{explicar} ao menos uma atribui\c{c}\~ao do cargo envolvido;
(iii) \textbf{distinguir} os pap\'eis do Legislativo e do Executivo; e (iv)
reconhecer que pol\'iticas p\'ublicas envolvem escolhas, limites e consequ\^encias.

O projeto se alinha \`a Agenda 2030 da ONU, dialogando sobretudo com o
\textbf{ODS~16} (Paz, Justi\c{c}a e Institui\c{c}\~oes Eficazes) e o \textbf{ODS~4}
(Educa\c{c}\~ao de Qualidade), al\'em dos ODS~11 e~17, partindo da premissa de que
cidad\~aos que compreendem a divis\~ao de compet\^encias fiscalizam melhor e participam
mais~\cite{ods16,ods4,agenda2030}. Princ\'ipios editoriais norteiam o conte\'udo:
personagens e falas fict\'icios, aus\^encia de partido ou pedido de voto, e revis\~ao
com fontes oficiais~\cite{constituicao,camara,senado}.

\section{Desenvolvimento}

\paragraph{Fundamenta\c{c}\~ao de \emph{game design}.} O design partiu do arcabou\c{c}o
\textbf{MDA} (Mec\^anicas $\rightarrow$ Din\^amicas $\rightarrow$ Est\'eticas), que
estrutura o jogo de baixo para cima enquanto o jogador o experimenta de cima para
baixo~\cite{mda}. As defini\c{c}\~oes de jogo, divers\~ao e experi\^encia, bem como as
``lentes'' de projeto, seguiram Schell~\cite{schell}, e as emo\c{c}\~oes-alvo foram
mapeadas pelas \emph{Four Fun Keys} de Lazzaro --- com \^enfase no \emph{Serious Fun},
a chave mais alinhada a jogos educativos~\cite{lazzaro}. A justificativa do
``jogo como prot\'otipo'' e as decis\~oes de \emph{game feel} apoiaram-se em
Fullerton~\cite{fullerton}, e o prop\'osito educativo foi fundamentado na literatura
de \emph{serious games}, que trata educa\c{c}\~ao e entretenimento como complementares,
tecendo a mensagem na experi\^encia em vez de preg\'a-la~\cite{seriousgames}.

\paragraph{Narrativa e personagens.} A constru\c{c}\~ao narrativa seguiu Nakamura,
distinguindo a narrativa \emph{no} jogo da que se forma \emph{ao redor} dele, e os
quatro modos de interatividade narrativa de Zimmerman~\cite{nakamura,zimmerman}.
A protagonista, Joana, foi projetada como \emph{avatar habit\'avel} e os NPCs como
pap\'eis sociais, conforme Isbister~\cite{isbister}. A quarta apresenta\c{c}\~ao conecta
a perspectiva do projetista a agentes racionais, estrat\'egicos, sociais e orientados
por valores parcialmente ordenados, com base na aula de modelagem de agentes de
Silva~\cite{silva}.

\paragraph{Implementa\c{c}\~ao t\'ecnica.} O prot\'otipo foi desenvolvido em
\textbf{Python~3} com a biblioteca \textbf{Arcade}~\cite{arcade} (janela,
renderiza\c{c}\~ao, \emph{sprites}, entrada e \'audio) e testes automatizados com
\textbf{pytest}~\cite{pytest}. O c\'odigo (\texttt{src/}) \'e modular: \texttt{config},
\texttt{assets}/\texttt{modular}/\texttt{citytiles}, \texttt{audio}, \texttt{rules},
\texttt{save\_manager}, \texttt{ui}, al\'em de \texttt{content/} (miss\~oes, cargos,
problemas, \emph{side quests}), \texttt{views/} (telas) e \texttt{world/} (interiores).
O progresso \'e persistido em JSON. Os \emph{assets} visuais utilizam pacotes
Kenney sob licen\c{c}a CC0~\cite{kenney,cc0}, e o jogo funciona sem \emph{assets}
externos de \'audio.

\paragraph{Mec\^anicas.} O jogo oferece movimenta\c{c}\~ao \emph{top-down}
(setas/\texttt{WASD}), di\'alogos e miss\~oes com NPCs que apresentam um problema e uma
pergunta de m\'ultipla escolha sobre o n\'ivel de governo respons\'avel, \emph{side
quests} com moradores, um \emph{Di\'ario de Cidadania} consult\'avel e um sistema de
provas (problema, testemunho e lei). A progress\~ao acompanha os tr\^es n\'iveis de
governo: Munic\'ipio, Estado e Uni\~ao.

\paragraph{Documenta\c{c}\~ao e apresenta\c{c}\~ao.} Paralelamente ao jogo, uma aplica\c{c}\~ao
web (React + Vite + Tailwind) re\'une quatro apresenta\c{c}\~oes --- prot\'otipo,
\emph{game design}, narrativa/personagens e vis\~ao conceitual/agentes --- que
explicitam as fontes, distinguem os conceitos das aulas de suas aplica\c{c}\~oes a
Aurora e identificam as propostas de evolu\c{c}\~ao ainda n\~ao implementadas.

\section{Resultados}

Entregou-se um \emph{vertical slice} funcional e execut\'avel, com a pra\c{c}a de
Aurora do Brasil, movimenta\c{c}\~ao \emph{top-down}, personagens fict\'icios, sistema de
di\'alogos e uma miss\~ao completa sobre a responsabilidade municipal por uma unidade
de sa\'ude, al\'em de \emph{side quests}, Di\'ario de Cidadania, salvamento em m\'ultiplos
\emph{slots} e configura\c{c}\~oes. A arquitetura modular separou conte\'udo educativo
(dados) da l\'ogica de jogo, facilitando a expans\~ao para as etapas estadual e federal.
As regras das miss\~oes s\~ao cobertas por testes automatizados, garantindo a
consist\^encia entre problema, n\'ivel de governo e resposta correta.

A aplica\c{c}\~ao de apresenta\c{c}\~ao consolidou a rastreabilidade entre teoria e projeto,
vinculando cada decis\~ao de design a um referencial acad\^emico e aos impactos
esperados no jogador --- tratados explicitamente como hip\'oteses a validar, com
cuidados de acessibilidade, fidelidade c\'ivica e preven\c{c}\~ao de vieses. Como
trabalhos futuros, prop\~oe-se completar as etapas estadual e federal, conduzir
\emph{playtesting} formal segundo as cinco perguntas de Schell~\cite{schell} e
explicitar, na jogabilidade, conflitos de valores parcialmente ordenados conforme a
modelagem de agentes~\cite{silva}. O impacto de aprendizagem e social permanece
condicionado a avalia\c{c}\~ao emp\'irica pr\'opria, n\~ao decorrendo automaticamente do jogo.

\balance
\printbibliography[title={Refer\^encias}]

\end{document}
